% 근거: analysis.json, chapter4-prep-v1/A5/*,
%       chapter4-prep-v1/A6/human-review-boundary.json
\subsection{자동 답변 품질 및 자원 사용 특성}
\label{sec:ch4_answer_resources}

Jev 조건에서는 최종 답변의 자동 평가 점수도 높았다
(표~\ref{tab:ch4_main_results} 하단). LF-Eval은 수정 CE 67.57점에서
Jev 70.15점으로 2.58점 증가하였으며, 차이의 탐색적 95\% 신뢰구간은
$[+0.37,\,+4.79]$점이었다.
Token F1은 46.20점에서 47.31점으로 1.12점 증가하였고,
차이의 탐색적 95\% 신뢰구간은 $[+0.02,\,+2.27]$점이었다.

Token F1은 정규화한 답변과 정답의 공백 토큰 중복을 측정하며,
생성 모델의 사용량 토큰과는 정의가 다르다. 공식 코드와 같은 800자
평가 규칙을 적용하였다. 이 지표의 centered-bootstrap 양측 $p$값은
약 0.0525였다. percentile 구간과 해당 검정은 동일 방법의 구간ㆍ검정
쌍이 아니므로, 자동 지표 전체를 통계적으로 유의한 개선으로 묶어
해석하지 않았다.

LF-Eval은 공식 정답 답변과 정답 근거(gold context)를 기준으로 한
평가이다. 실제 제공한 근거가 답변의 주장을 지지하는지는 별도의
사람 평가가 필요하다. 기존 60문항ㆍ세 조건의 총 180답변에 대한
블라인드 양식은 준비되었으나 실제 판단은 아직 없다. 따라서 사람
정확도ㆍ근거 지지율ㆍ평가자 일치도는 미측정 상태이며, 본 절의
품질 결과는 완료된 자동 평가에 한정된다.

자동 점수의 증가와 함께 평균 문맥ㆍ입력량도 소폭 늘었다.
문맥 토큰의 평균은 수정 CE 203.67개, Jev 215.60개였으며,
중앙값은 각각 164개와 170.5개, 사분위 범위(IQR)는
133.75개와 142.50개였다. 문항별 Jev$-$수정 CE의 대응 차이는
평균 $+11.93$개, 중앙값 0개, IQR 2.50개였다.
질문ㆍ지시문ㆍ예시까지 포함한 전체 API 입력 토큰은 평균
2302.67개와 2314.58개, 출력 토큰은 72.16개와 73.98개였다.
평균 입력량의 절감은 관찰되지 않았으며, 이 기술통계는 문맥 길이가
품질 차이의 원인임을 보여주는 결과도 아니다.

두 조건의 선택 목록 길이는 문항별로 같았으며, 각 조건의 합계는
478개였다. 문항 내 고유 근거 수의 합계는 각각 407개와 409개,
중복 출현은 71회와 69회였다. 선택 개수를 줄이는 Filtering을 적용한
결과가 아니라 선택 내용과 중복 양상의 차이이다.
문맥 토큰은 실제 제공 목록의 순서와 중복을 유지하고 줄바꿈으로
연결한 문자열에 \texttt{o200k\_base} 토크나이저를 적용해 산출하였다.
전체 입력ㆍ출력량은 생성 응답의 사용량 기록을 사용하였다.

이 사용량은 생성 단계의 입출력 요구량으로, 신규 청구액과 구별된다.
생성 응답 캐시를 재사용한 문항은 수정 CE 72개, Jev 122개였으며,
캐시 행도 동일 입력이 요구하는 양을 기술하는 데 포함하였다.
선택ㆍ점수화ㆍ평가ㆍ재시도와 로컬 CE의 운영 단가를 모두 반영한
독립 서비스 비용은 산출하지 않았다. 캐시와 요청 제한 대기가 포함된
이번 품질 실행의 시간도 공정한 온라인 지연 비교에 사용할 수 없다.
